\uuid{PhQs}
\exo7id{6890}
\titre{exo7 6890}
\auteur{ruette}
\organisation{exo7}
\datecreate{2013-01-24}
\isIndication{false}
\isCorrection{false}
\chapitre{Probabilité discrète}
\sousChapitre{Probabilité et dénombrement}
\cpge{mpsi-mp2i.probabilites-univers-fini.b, pcsi.probabilites-variables-aleatoires.b, ptsi.probabilites-variables-aleatoires.b, mpsi-mp2i.probabilites-univers-fini.c, pcsi.probabilites-variables-aleatoires.c, ptsi.probabilites-variables-aleatoires.c, bcpst1.concepts-base-probabilites.a, bcpst1.concepts-base-probabilites.b, mp-mpi.variables-aleatoires-discretes.c, pc.probabilites-lois-usuelles.c, psi.probabilites-lois-usuelles.c, pt.probabilites-lois-usuelles.d}
\module{Probabilité et statistique}
\niveau{L2}
\difficulte{}

\contenu{
\texte{
Au loto, on choisit 6 numéros principaux, qui sont 6 nombres différents 
entre 1 et 49, et un numéro 
complémentaire, qui est un nombre entre 1 et 49 différent des 6 précédents. 
Quelle est la probabilité d'avoir :
}
\begin{enumerate}
    \item \question{les six bons numéros principaux ?}
    \item \question{cinq bons numéros parmi les 6 principaux?}
    \item \question{cinq bons numéros parmi les principaux 
et le bon numéro complémentaire ?}
\end{enumerate}
}
